\documentclass{standalone}
\usepackage{circuitikz}
\usetikzlibrary{calc}
\definecolor{circuitnotemark}{HTML}{FE00FE}
\makeatletter
\pgfdeclareshape{smacoax}{
  \anchor{center}{\pgfpointorigin}
  \anchor{text}{\pgfpointorigin}
  \anchor{north}{\pgfpoint{0cm}{0.3cm}}
  \anchor{south}{\pgfpoint{0cm}{-0.3cm}}
  \anchor{east}{\pgfpoint{0.3cm}{0cm}}
  \anchor{west}{\pgfpoint{-0.3cm}{0cm}}
  \anchor{pin 1}{\pgfpoint{-0.7cm}{0cm}}
  \anchor{pin 2}{\pgfpoint{0cm}{-0.7cm}}
  \anchor{bpin 1}{\pgfpoint{-0.3cm}{0cm}}
  \anchor{bpin 2}{\pgfpoint{0cm}{-0.3cm}}
  \backgroundpath{
    \pgfpathmoveto{\pgfpoint{-0.274cm}{-0.122cm}}
    \pgfpatharc{204}{516}{0.3cm}
    \pgfpathmoveto{\pgfpoint{-0.7cm}{0cm}}\pgfpathlineto{\pgfpointorigin}
    \pgfpathmoveto{\pgfpoint{0cm}{-0.7cm}}\pgfpathlineto{\pgfpoint{0cm}{-0.3cm}}
  }
  \foregroundpath{
    \pgfpathcircle{\pgfpointorigin}{0.07cm}
    \pgfusepath{fill}
  }
}
\makeatother
\begin{document}
\begin{circuitikz}[american, line width=0.8pt]
\ctikzset{bipoles/length=1.2cm}
\ctikzset{grounds/scale=1.36}
\coordinate (x2y3) at (2,-4);
\coordinate (x4y3) at (6,-4);
\coordinate (x4y5) at (6,-8);
\coordinate (x5y3) at (8,-4);
\coordinate (x7y3) at (12,-4);
\coordinate (x8y3) at (14,-4);
\coordinate (x8y5) at (14,-8);
\coordinate (x10y3) at (18,-4);
\coordinate (x2y4) at (2,-6);
\coordinate (x10y4) at (18,-6);
\node[smacoax, draw, xscale=-1] (part-J1) at (x2y3) {}; % line 3
\node[font=\scriptsize, anchor=south] at (part-J1.north) {$\mathrm{CH0}$}; % line 3
\node[anchor=east] at (part-J1.east) {$J_{1}$}; % line 3
\draw (x4y3) to[C, l_=$C_{1}$, a^=$47\,\mathrm{p}\mathrm{F}$] (x4y5); % line 4
\draw (x5y3) to[L, l_=$L_{1}$, a^=$235\,\mathrm{n}\mathrm{H}$] (x7y3); % line 5
\draw (x8y3) to[C, l_=$C_{2}$, a^=$47\,\mathrm{p}\mathrm{F}$] (x8y5); % line 6
\node[smacoax, draw] (part-J2) at (x10y3) {}; % line 7
\node[font=\scriptsize, anchor=south] at (part-J2.north) {$\mathrm{CH1}$}; % line 7
\node[anchor=west] at (part-J2.east) {$J_{2}$}; % line 7
\node[ground] at (x2y4) {}; % line 8
\node[ground] at (x10y4) {}; % line 9
\node[ground] at (x4y5) {}; % line 10
\node[ground] at (x8y5) {}; % line 11
\draw (part-J1.pin 1) -- (x4y3); % line 13
\draw (x4y3) -- (x5y3); % line 13
\draw (x7y3) -- (x8y3); % line 14
\draw (x8y3) -- (part-J2.pin 1); % line 14
\draw (part-J1.pin 2) -- (x2y4); % line 15
\draw (part-J2.pin 2) -- (x10y4); % line 16
\node[circ] at (x4y3) {};
\node[circ] at (x8y3) {};
\node[anchor=south west, inner sep=2pt, yshift=4pt, circuitnotemark, font=\LARGE\bfseries] (circuittitle) at (current bounding box.north west) {X};
\path (circuittitle.south west) rectangle ++(15.453,0.853);
\node[anchor=north east, inner sep=2pt, circuitnotemark, font=\large] (circuitstamp) at ([yshift=-0.593cm]current bounding box.south east) {X};
\path (circuitstamp.north east) rectangle ++(-5.08,-0.593);
\end{circuitikz}
\end{document}
